\section{State Adoption Pathways}
\label{sec:state}

In the absence of federal legislation, states may adopt algorithmic redistricting through three mechanisms: ordinary legislation, commission action, and court-ordered remedies. The availability of each mechanism depends on the state's constitutional structure.

\subsection{Constitutional Typology}

Fifty state constitutions fall into five categories based on their compatibility with algorithmic redistricting, which we analyze in turn.

\textit{Type I: Permissive (no obstacles).} Twenty-eight states have constitutional provisions that are either silent on redistricting method or grant the legislature plenary authority over congressional redistricting. Examples include Texas, Georgia, Florida (congressional only), and Wisconsin. In these states, the legislature may mandate algorithmic redistricting by ordinary statute. No constitutional amendment is required; no commission need be consulted. The legal analysis is uncomplicated: a rational basis is supplied by the legislature's interest in reducing litigation, increasing public trust, and demonstrating impartiality. Model statute text provided in Section~\ref{sec:model} can be adopted with minimal state-specific modification.

\textit{Type II: Criteria-based (explicit constitutional standards).} Fourteen states have constitutions that enumerate redistricting criteria---compactness, contiguity, preserving political subdivisions, or communities of interest---without specifying the process by which those criteria must be implemented. Examples include Colorado (Art. V, \S\,44, as amended by Amendments Y and Z), Michigan (Art. IV, \S\,6, as amended by Prop 2), and Virginia (Art. II, \S\,6a). In these states, the argument is stronger: the algorithmic approach \textit{satisfies} the constitutional criteria by design. A compactness-optimizing algorithm that preserves contiguity and respects census-tract boundaries satisfies these constitutional mandates. Adoption requires a statute specifying that the algorithm is the state's chosen method of satisfying its constitutional redistricting obligations.

\textit{Type III: Commission-mandated (constitutional commission requirement).} Six states require redistricting by an independent commission, but do not specify the commission's internal methodology: Arizona, California, Idaho, Montana (as of 2022), Oregon, and Washington. In these states, the commission retains full authority over methodology. Each commission may adopt algorithmic redistricting by internal resolution, specifying that the certified algorithm produces the baseline map from which the commission may make adjustments pursuant to its constitutional criteria. The Arizona Independent Redistricting Commission's adoption process would be a model: commissioners vote to certify the algorithm, promulgate rules governing the adjustment period, and conduct public hearings on proposed adjustments before adopting the final map.\footnote{See \textit{Arizona State Legislature v.\ Arizona Independent Redistricting Commission}, 576 U.S. 787 (2015) (upholding commission's broad authority over redistricting methodology).}

\textit{Type IV: Hybrid (commission plus legislative ratification).} Four states have constitutional structures requiring commission action followed by legislative approval: New York (Art. III, \S\S\,4-5, as amended by 2021 referendum), New Jersey (Art. IV, \S\,3), Montana (Art. V, \S\,14, pre-2022), and Hawaii (Art. IV, \S\,2). In these states, the commission would adopt an algorithmic baseline and present it to the legislature for ratification. The legislature would be constrained in its ability to modify the algorithmic output by the commission's constitutional mandate to produce compact, contiguous districts.

\textit{Type V: Restrictive (explicit process requirements).} A small number of states have constitutional language that may require human deliberation in redistricting. Maine's constitution requires that the redistricting commission's decisions be reached through a deliberative process with specified majority requirements. No state constitution explicitly prohibits algorithmic methods, but some language arguably requires decision-making by named bodies in ways that a fully automated process might not satisfy. In these states, algorithmic redistricting should be framed as the commission's tool of choice rather than a replacement for commission decision-making: the commission uses the algorithm to generate a baseline, deliberates publicly over any adjustments, and adopts the final map by the constitutionally required majority.

\subsection{State Constitutional Amendment Strategy}

For the eight states where the existing constitutional text creates meaningful obstacles---typically states with prescriptive commission structures that describe an inherently deliberative, iterative process---a constitutional amendment is the cleanest solution. The amendment need not be elaborate. A single provision authorizing the legislature or commission to adopt algorithmic redistricting methods, subject to the state's existing substantive criteria, is sufficient.

The procedural amendment language (to be added to existing redistricting provisions) would read: ``The [Legislature / Commission] may adopt computer-based algorithmic methods for producing congressional and state legislative district maps, provided that such methods use only population and geographic data from federal census sources, do not incorporate data on voters' partisan affiliation or voting history, and produce districts that satisfy all substantive requirements of this Article.'' Transparency and public auditing requirements should be added to ensure that the amendment does not create new opacity.

\subsection{Court-Ordered Algorithmic Redistricting}

The most immediate pathway to algorithmic redistricting in the near term may be court-ordered remedial adoption. When a court strikes down a congressional district plan as violating a state constitution---as Pennsylvania's Supreme Court did in 2018, North Carolina's in 2022, and New York's appellate division in 2021---it must specify a remedy. The existing remedy practice is for the court to appoint a special master who draws a remedial map using criteria specified by the court.

Algorithmic redistricting is an ideal remedial tool for three reasons. First, it eliminates the special master's own subjectivity. Special masters, however experienced and well-intentioned, make line-drawing choices that parties will challenge as reflecting their own preferences or biases. An algorithm's choices are fully documented in its code and constrained by its mathematical formulation, reducing the scope for such challenges. Second, it is faster. Running the algorithm on census data produces a complete map within hours. Special master processes that involve map-drawing, revision, and response to objections typically take months. Third, it produces maps that no court need defend on normative grounds---the court can simply note that the algorithm optimizes for the criteria its own precedents have identified as constitutionally required, and that its outputs satisfy one-person-one-vote requirements.

The model court order in the model statute section specifies the algorithmic method, requires the relevant state agency to run the certified algorithm within 14 days, provides for a 21-day public comment period on the output, and permits parties to file objections limited to population-balance or geographic-adjacency errors. Courts retain jurisdiction to resolve objections and to order reversion to the unmodified algorithmic output if proposed modifications cannot be agreed.
